\def\titlestr{Quantenmechanik}
\ifx\makebook\undefined
	% ----------------------------------------------------------------
% ----------------------------------------------------------------
% AMS-LaTeX Paper ************************************************f
% **** -----------------------------------------------------------
% documentclass[fleqn]{report}
%\RequirePackage{silence}
%\WarningFilter{lualatex}{Token not allowed}
% Einzelne Warnung eines Pakets unterdrücken:
%\WarningFilter{pdftex}{destination with}

% Alle Warnungen eines bestimmten Pakets unterdrücken:
%\WarningsOff[lualatex]
\PassOptionsToPackage{unicode}{hyperref}
%\PassOptionsToPackage{naturalnames}{hyperref}
\documentclass[a4paper,fleqn]{article}
%\documentclass[a4paper,fleqn]{scrreprt}
%\documentclass{styles/wissdoc-kw-eng}
\setlength{\oddsidemargin}{0.0cm}
\setlength{\evensidemargin}{0.0cm}
\setlength{\paperheight}{297mm}
%\usepackage[utf8]{inputenc}
\usepackage{german}
\usepackage{graphicx} 
\usepackage{amsmath}
\usepackage{amssymb}
\usepackage{alltt}
\usepackage{color}
\usepackage{fancyvrb}
\usepackage{verbatim}
\usepackage{mathtools}
\usepackage[compat=1.1.0]{tikz-feynman}
\DeclarePairedDelimiter\bra{\langle}{\rvert}
\DeclarePairedDelimiter\ket{\lvert}{\rangle}
\DeclarePairedDelimiterX\braket[2]{\langle}{\rangle}{#1\,\delimsize\vert\,\mathopen{}#2}
\DeclarePairedDelimiterX\expect[1]{\langle}{\rangle}{#1}
%\newcounter{thm}[section]
%\newtheorem{theorem}[thm]{thm}
%\newtheorem{lemma}[theorem]{Lemma}
%\newtheorem{proposition}[theorem]{Hilfssatz}
%\newtheorem{corollary}[theorem]{Corollary}
%\newtheorem{definition}[theorem]{Definition}
%\newenvironment{theorem}[Theorem]{\begin{trivlist}
%\item[\hskip \labelsep {\textit \textit{#1}}]}{\end{trivlist}}
%\newenvironment{proof}[1][Beweis]{\begin{trivlist}
%\item[\hskip \labelsep {\bfseries #1}]}{\end{trivlist}}
%\item[\hskip \labelsep {\bfseries #1}]}{\end{trivlist}}
\newcommand*\codeRed[1]{\textcolor[rgb]{1,0,0}{\textbf{#1}}}
\newcommand*\codeGreen[1]{\textcolor[rgb]{0,0.5,0}{\textbf{#1}}}
\newcommand{\lb}{[}
\newcommand{\rb}{]}
\newcommand{\codeCall}[1]{\textbf{\texttt{#1}}}
\newcommand{\mcL}{\mathcal{L}}
\newcommand{\kb}{k_{\mathrm{B}}}
%\newcommand{\titlemath}[1]{\texorpdfstring{#1}{}}
\newcommand{\titlemath}[2]{\texorpdfstring{#1}{#2}}
\newcommand{\pmat}[1]{\begin{pmatrix}#1\end{pmatrix}}


%\newenvironment{\vspace{1ex}\noindent{\bf Proof}\hspace{0.5em}}
%	{\hfill\qed\vspace{1ex}}
 

\newenvironment{example}[1][Example]{\begin{trivlist}
\item[\hskip \labelsep {\bfseries #1}]}{\end{trivlist}}
\newenvironment{remark}[1][Remark]{\begin{trivlist}
\item[\hskip \labelsep {\bfseries #1}]}{\end{trivlist}}

%\newcommand{\qed}{\nobreak \ifvmode \relax \else
%      \ifdim\lastskip<1.5em \hskip-\lastskip
%      \hskip1.5em plus0em minus0.5em \fi \nobreak
%      \vrule height0.75em width0.5em depth0.25em\fi}
%\newcommand{\qed}{$\square$}
\newcommand{\qt}{\texttt}

\newcommand{\bibdef}[4]{\bibitem{#1} {\bf #2:} {\it #3}\\#4}


\newcommand{\gauss}[1]{\lfloor #1\rfloor}
\numberwithin{equation}{section}


%\setcounter{equation}{1}

\setcounter{tocdepth}{4}
\setcounter{secnumdepth}{4}

\newcommand{\ltri}{\vartriangleleft}
\usepackage{amsthm}
\newtheorem{theorem}{Theorem}[section]
\newtheorem{definition}{Definition}[section]
\renewcommand{\proofname}{Beweis}
\definecolor{dblue}{RGB}{150,50,150}
%\usepackage[colorlinks=true, linkcolor=dblue, urlcolor=cyan]{hyperref}
\usepackage{hyperref}
%\usepackage{cleveref}

%\usepackage{crossreftools}

%\addto\captionsngerman{%
 \renewcommand{\abstractname}{Vorbemerkung}
%}
	\begin{document}
	\title{\titlestr}
	\author{Hugo Meder}
	\maketitle
	\tableofcontents
	\newpage
	%\begin{abstract}
	Dieser Text ist im Rahmen meiner Vortragsreihe 'Ästhetik der Rationalität' entstanden. Es handelt sich \textbf{nicht} um den Text des Vortrags selbst. Er enthält aber die Mathematik und die Physik die im Vortrag dargestellt wird. Ziel ist ein doppeltes:
\begin{itemize}
	\item Dem Interessierten Zugang zu einer formal exakten Formulierung des Inhalts des Vortrags zu gewähren.
	\item Für mich selbst erfüllt das Schreiben den Zweck mir die Struktur und die logischen Zusammenhänge klar zu machen. Diese Struktur führt dann auch dazu Reihenfolge und Darstellung des Vortrags zu gestalten.
\end{itemize}

	%\end{abstract}
	\newpage

\else
	\begin{chapter}{\titlestr}
\fi
\begin{section}{Plancksche Hohlraumstrahlung}
\begin{subsection}{Elektromagnetische Schwingungsmodi in einem Volumen}
Wir nehmen zunächst an in einem Vakuum in einem Hohlraum schwingen elektromagnetische Wellen im thermischen Gleichgewicht der Temperatur $T$. Das beschränkte Volumen $V$ führt zu Randbedingungen. Diese
haben wiederum die Folge, dass nur bestimmte Schwingungsmodi erlaubt sind. Das entsprechende Randwertproblem ist schwierig zu lösen, selbst für so einfache Geometrien wie einem Quader.

Es gibt aber die zyklische Randbedingung. In einer Dimension bedeutet das
\begin{eqnarray*}
E(0,t) &=& E(l,t)
\end{eqnarray*}
und in drei Dimensionen:
\begin{eqnarray*}
E(0,y,z,t) &=& E(l_x,y,z,t)\\
E(x,0,z,t) &=& E(x,l_y,z,t)\\
E(x,y,0,t) &=& E(x,y,z,-l_zt)
\end{eqnarray*}
Eine  ebene Welle siht so aus:
\begin{eqnarray*}
E(x,y,z,t) &=& E\sin(xk_x+yk_y+zk_z-\omega t+\phi)
\end{eqnarray*}
Die zyklische Randbedingung führt nun zu
\begin{eqnarray*}
\sin(yk_y+zk_z-\omega t+\phi) &=& \sin(l_xk_x+yk_y+zk_z-\omega t+\phi)\\
\sin(xk_x+zk_z-\omega t+\phi) &=& \sin(xk_x+l_yk_y+zk_z-\omega t+\phi)\\
\sin(xk_x+yk_y-\omega t+\phi) &=& \sin(xk_x+yk_y+l_zk_z-\omega t+\phi)\\
k_xl_x &=& n_x2\pi\\
k_yl_y &=& n_y2\pi\\
k_zl_z &=& n_z2\pi\\
k_x &=& n_x\frac{2\pi}{l_x}\\
k_y &=& n_y\frac{2\pi}{l_y}\\
k_z &=& n_z\frac{2\pi}{l_z}
\end{eqnarray*}
Die Dichte der Modi ist also
\begin{eqnarray*}
\rho(\vec{k}) &=& \frac{l_xl_yl_z}{(2\pi)^3}
\end{eqnarray*}
Da elektromagnetisch Wellen zwei Polarisationsrichtungen haben, muss mit 2 multipliziert werden:
\begin{eqnarray*}
\rho(\vec{k}) &=& 2\frac{l_xl_yl_z}{(2\pi)^3}=2\frac{V}{(2\pi)^3}
\end{eqnarray*}
Die Anzahl der Modi in einer Kugelschale
zwischen $|\vec{k}|=\kappa$ und $|\vec{k}|=\kappa +\Delta\kappa$ ist
\begin{eqnarray*}
d N &=& \frac{4\pi}{3}\rho \kappa^2 d\kappa=\frac{4\pi}{3}2\frac{V}{(2\pi)^3} \kappa^2 d\kappa
=\frac{4}{3}\frac{V}{(2\pi)^2} \kappa^2 d\kappa
\end{eqnarray*}
Für jeden $\vec{k}$-Vektor ist die zugehörige Kreisfrequenz $\omega = c|\vec{k}|$. Daher ist die Anzahl der Modi mit einer Kreisfrequenz zwischen $\omega$ und $\omega+\delta \omega$:
\begin{eqnarray}
\label{state_density}
d N(\omega,V) &=& \frac{4}{3}\frac{V}{(2\pi)^2c^3} \omega^2 d\omega
\end{eqnarray}
Wir haben also $N$ Schwingungs-Modi mit der Kreisfrequenz $\omega$.
\end{subsection}

\begin{subsection}{Harmonischer Oszillator im Wärmebad}
Wir betrachten einen harmonischen Oszillator mit Masse $m$ und Federkonstante $f$. Die entsprechende Hamilton-Funktion ist
\begin{eqnarray*}
H(p,q) &=& \frac{1}{2}\frac{p^2}{m}+\frac{1}{2}fq^2
\end{eqnarray*}
und daraus ergeben sich die Bewegungsgleichungen
\begin{eqnarray*}
\dot{p} &=& -\frac{\partial H}{\partial q} = -fq\\
\dot{q} &=& \frac{\partial H}{\partial p} = \frac{p}{m}
\end{eqnarray*}
Und daraus die Kreisfrequenz:
\begin{eqnarray*}
\ddot{p} &=& -\frac{f}{m}p\\
q(t) &=& a\sin(\omega t)\\
p(t) &=& ma\omega\cos(\omega t)\\
\omega &=& \sqrt{\frac{f}{m}}
\end{eqnarray*}
$a$ ist die Amplitude der Schwingung. Die Energie ist
\begin{eqnarray*}
E(a) &=& H(p,q) = \frac{1}{2}\frac{\left(ma\omega\cos(\omega t)\right)^2}{m}+\frac{1}{2}f\left(a\sin(\omega t)\right)^2\\
&=& \frac{1}{2}a^2\left(\omega^2m^2\frac{\cos^2(\omega t)}{m}+f\sin^2(\omega t)\right)\\
&=& \frac{1}{2}a^2\left(\frac{f}{m}m^2\frac{\cos^2(\omega t)}{m}+f\sin^2(\omega t)\right)\\
&=& \frac{1}{2}a^2f\left(\cos^2(\omega t)+\sin^2(\omega t)\right)\\
&=& \frac{1}{2}a^2f
\end{eqnarray*}
Der von der Kurve $p(t)$ und $q(t)$ umschlossene Phasenraum $V$ (eine Ellipse) ist:
\begin{eqnarray*}
V &=& \pi ma^2\omega\\
a^2(E) &=& \frac{2E}{f}\\
V &=& 2\pi E\frac{m}{f}\omega=\frac{2\pi E}{\omega}=E\tau
\end{eqnarray*}
Wobei $\tau$ die Zeit für einen Schwingungs-Zyklus ist. Die Phasenraumdichte $\omega(E)=\frac{dV}{dE}$ hängt nicht von der Energie ab:
\begin{eqnarray*}
\omega(E)&=&\frac{2\pi}{\omega}=\tau
\end{eqnarray*}
Die Wahrscheinlichkeits-Dichte im Wärmebad ist allgemein
\begin{eqnarray*}
\sigma(E,\beta)&=& \frac{\omega(E)e^{-\beta E}}{Z(\beta)}\\
Z(\beta) &=& \int\omega(E)e^{-\beta E}dE
\end{eqnarray*}
und für den Oszillator
\begin{eqnarray*}
\sigma(E)&=& \frac{\tau e^{-\beta E}}{Z(\beta)}\\
Z(\beta) &=& \int_0^\infty\tau e^{-\beta E}dE=\left(-\frac{\tau}{\beta}e^{-\beta E}\right|_0^\infty=\frac{\tau}{\beta}\\
\sigma(E)&=& \beta e^{-\beta E}
\end{eqnarray*}
Der Erwartungswert für die Energie ist somit
\begin{eqnarray*}
\expect*{E(\omega,T)} &=& \int_0^\infty E \sigma(E) dE=\int_0^\infty \beta Ee^{-\beta E} dE=\frac{1}{\beta} = \kb T
\end{eqnarray*}
\end{subsection}

\begin{subsection}{Die Ultraviolett-Katastrophe}
Mit der Zustands-Dichte \ref{state_density}
\begin{eqnarray*}
d N(\omega,V) &=& \frac{4}{3}\frac{V}{(2\pi)^2c^3} \omega^2 d\omega
\end{eqnarray*}
ergibt sich eine Energie-Dichte
\begin{eqnarray*}
d E(\omega,V,T) &=& \frac{4}{3}\frac{V}{(2\pi)^2c^3} \omega^2 \expect*{E(\omega,T)} d\omega\\
&=& \frac{4}{3}\frac{V}{(2\pi)^2c^3} \omega^2 \kb T d\omega
\end{eqnarray*}
Die Gesamtenergie $E(V,)$ ist
\begin{eqnarray*}
E(V,T) &=& \int_0^\infty d E(\omega,V,T) )=\int_0^\infty\frac{4}{3}\frac{V}{(2\pi)^2c^3} \omega^2 \kb T d\omega\\
&=&\frac{4}{3}\frac{V}{(2\pi)^2c^3}\kb T \int_0^\infty\omega^2 d\omega =\infty
\end{eqnarray*}
Dies ist die sogenannte Ultraviolett-Katastrophe.
\end{subsection}

\begin{subsection}{Quantisierter harmonischer Oszillator}
Auf der Suche nach einer Lösung hat Max Planck den Fall untersucht, bei dem ein Harmonischer Oszillator nur diskrete Energien annehmen kann:
\begin{eqnarray*}
E(\omega,n) &=& n\Delta E(\omega)
\end{eqnarray*}
Die Wahrscheinlichkeits-Verteilung ist dann entsprechend auch diskret:
\begin{eqnarray*}
p(\omega,n,\beta)&=& \frac{e^{-\beta E(\omega,n)}}{Z(\beta,\omega)}\\
Z(\omega,\beta) &=& \sum_{n=0}^\infty e^{-\beta E(\omega,n)}= \sum_{n=0}^\infty e^{-\beta\left(n\Delta E(\omega)\right)}\\
&=& e^{-\beta E_0(\omega)}\sum_{n=0}^\infty e^{-\beta n\Delta E(\omega)}
\end{eqnarray*}
Mit
\begin{eqnarray*}
x &=& e^{-\beta \Delta E(\omega)}
\end{eqnarray*}
wird
\begin{eqnarray*}
Z(\omega,\beta) &=& e^{-\beta E_0(\omega)}\sum_{n=0}^\infty x^n\\
\sigma &=& \sum_{n=0}^\infty x^n\\
x\sigma &=& \sum_{n=1}^\infty x^n\\
(1-x)\sigma &=& 1\\
\sigma &=& \frac{1}{1-x}\\
Z(\omega,\beta) &=& \frac{e^{-\beta E_0(\omega)}}{1-e^{-\beta \Delta E(\omega)}}\\
p(\omega,n,\beta)&=& \frac{e^{-\beta E(\omega,n)}}{Z(\beta,\omega)}
=e^{-\beta n\Delta E(\omega)}\left(1-e^{-\beta \Delta E(\omega)}\right)
\end{eqnarray*}
Und der Erwartungswert $\expect{E(\omega,T)}$ wird
\begin{eqnarray*}
\expect{E(\omega,T)} &=& \sum_{n=0}^\infty \left(n\Delta E(\omega)\right)p(\omega,n,\beta)\\
&=& \sum_{n=0}^\infty n\Delta E(\omega)p(\omega,n,\beta)\\
&=& \left(1-e^{-\beta \Delta E(\omega)}\right)\Delta E(\omega)\sum_{n=0}^\infty n e^{-\beta n\Delta E(\omega)}\\
&=& \left(1-e^{-\beta \Delta E(\omega)}\right)\Delta E(\omega)\sum_{n=0}^\infty n x^n\\
\frac{d}{dx}\sum_{n=0}^\infty x^n &=& \sum_{n=0}^\infty nx^{n-1}\\
\sum_{n=0}^\infty n x^n &=& x\frac{d}{dx}\sum_{n=0}^\infty x^n
=x\frac{d}{dx}\frac{1}{1-x}=\frac{x}{(1-x)^2}\\
&=&\frac{e^{-\beta \Delta E(\omega)}}{\left(1-e^{-\beta \Delta E(\omega)}\right)^2}\\
\expect{E(\omega,T)} &=& \frac{e^{-\beta \Delta E(\omega)}}{1-e^{-\beta \Delta E(\omega)}}\Delta E(\omega)\\
&=& \frac{1}{e^{\beta \Delta E(\omega)}-1}\Delta E(\omega)
\end{eqnarray*}
Ist $\beta \Delta E(\omega)\ll 1$ also $\Delta E(\omega)\ll \kb T$ wird
\begin{eqnarray*}
\expect{E(\omega,T)} &=& \kb T
\end{eqnarray*}
Für $\beta \Delta E(\omega)\gg 1$ ergibt sich
\begin{eqnarray*}
\expect{E(\omega,T)} &=& 0
\end{eqnarray*}
\end{subsection}

\begin{subsection}{Die Plancksche Formel}
Zusammen mit \ref{state_density}
\begin{eqnarray*}
d N(\omega,V) &=& \frac{4}{3}\frac{V}{(2\pi)^2c^3} \omega^2 d\omega
\end{eqnarray*}
ergibt sich eine Energie-Dichte
\begin{eqnarray*}
d E(\omega,V,T) &=& \frac{4}{3}\frac{V}{(2\pi)^2c^3} \omega^2 \expect*{E(\omega,T)} d\omega\\
&=& \frac{4}{3}\frac{V}{(2\pi)^2c^3} \omega^2 \frac{1}{e^{\beta \Delta E(\omega)}-1}\Delta E(\omega) d\omega
\end{eqnarray*}
Ein Vergleich mit experimentellen Daten legt
\begin{eqnarray*}
\Delta E(\omega) &=& \frac{h\omega}{2\pi} = h\nu
\end{eqnarray*}
nahe, wobei $\nu$ die Frequenz ist.
\begin{eqnarray*}
d E(\nu,V,T) &=& \frac{4}{3}\frac{V}{(2\pi)^2c^3} (2\pi)^2\nu^2 \frac{1}{e^{\beta h\nu-1}}h\nu 2\pi d\nu\\
&=& \frac{8h\pi}{3}\frac{V}{c^3} \nu^3 \frac{1}{e^{\beta h\nu}-1}d\nu
\end{eqnarray*}
Bestimmt man $h$ aus den experimentellen Daten, dann stimmt diese Energie-Verteilung sehr gut mit den Experimenten überein.
$h$ heißt plancksches Wirkungs-Quantum. Auch der Vorfaktor stimmt, was nachträglich die zyklische Randbedingung stützt.
Anstelle der Frequenz wird auch gerne die Wellenlänge $\lambda=\frac{c}{\nu}$ verwendet:
\begin{eqnarray}
\label{lambda_density}
d E(\lambda,V,T) &=& \frac{8h\pi}{3}\frac{V}{c^3} \left(\frac{c}{\lambda}\right)^3 \frac{1}{e^{\frac{c\beta h}{\lambda}}-1}
\left|\frac{d\nu}{d\lambda}\right|d\lambda\cr
&=&\frac{8h\pi}{3}\frac{V}{\lambda^3} \frac{1}{e^{\frac{c\beta h}{\lambda}}-1}
\frac{c}{\lambda^2}d\lambda\cr
&=&\frac{8h\pi}{3}\frac{Vc}{\lambda^5} \frac{1}{e^{\frac{c\beta h}{\lambda}}-1}d\lambda
\end{eqnarray}
\end{subsection}

\begin{subsection}{Strahlungsintensität}
Wir wählen ein Volumenelement $d\mu$ im $\vec{k}$-Raum
\begin{eqnarray*}
d\mu &=& d\Omega |\vec{k}|^2d|\vec{k}|
\end{eqnarray*}
$d\Omega$ ist ein Oberflächenelement der Einheitskugel in $\vec{n}$-Richtung. Da $|\vec{k}|=\frac{\omega}{c}$ ist dies dasselbe wie
\begin{eqnarray*}
d\mu &=& d\Omega \frac{\omega^2}{c^3}d\omega
\end{eqnarray*}
Bei einer Zustandsdichte
\begin{eqnarray*}
\rho &=& 2\frac{V}{(2\pi)^3}
\end{eqnarray*}
gibt es
\begin{eqnarray*}
N &=& \rho d\mu = 2\frac{V}{(2\pi)^3}d\Omega \frac{\omega^2}{c^3}d\omega
=2\frac{V\omega^2}{(2\pi c)^3}d\Omega d\omega
\end{eqnarray*}
Modi in $\vec{n}$-Richtung. Die 'Photonen' bewegen sich mit der Lichtgeschwindigkeit $c\vec{n}$. Das Volumen, das bei dieser Geschwindigkeit und einem Zeitintervall $\Delta t$ ein Flächenelement $\vec{A}$ durchströmt ist
\begin{eqnarray*}
V &=& c\vec{n}\cdot\vec{A}\Delta t
\end{eqnarray*}
Und daher
\begin{eqnarray*}
N &=& \rho d\mu = 2\frac{c\vec{n}\cdot\vec{A}\Delta t\omega^2}{(2\pi c)^3}d\Omega d\omega\\
&=& 2\frac{\omega^2}{(2\pi )^3c^2}\vec{n}\cdot\vec{A}d\Omega d\omega\Delta t
\end{eqnarray*}
Und schließlich die Energie, die durch $\vec{A}$ fließt:
\begin{eqnarray*}
E(\omega,\vec{A},\Omega,\vec{n},T) &=& 2\frac{\omega^2}{(2\pi )^3c^2}\expect{E(\omega,T)}\vec{n}\cdot\vec{A}d\Omega d\omega\Delta t
\end{eqnarray*}
Mit $\nu=\frac{\omega}{2\pi}$, $\Omega=2\pi\nu$:
\begin{eqnarray*}
E(\nu,\vec{A},\Omega,\vec{n},T) &=& 2\frac{\nu^2}{c^2}\expect{E(\nu,T)}\vec{n}\cdot\vec{A}d\Omega d\nu\Delta t\\
\end{eqnarray*}
und
\begin{eqnarray*}
\expect{E(\nu,T)} &=& \frac{h\nu}{e^{\beta h\nu}-1}
\end{eqnarray*}
wird
\begin{eqnarray*}
E(\nu,\vec{A},\Omega,\vec{n},T) &=& 2\frac{h\nu^3}{c^2}\frac{1}{e^\frac{h\nu}{\kb T}-1}\vec{n}\cdot\vec{A}d\Omega d\nu\Delta t\\
P(\nu,\vec{A},\Omega,\vec{n},T) &=& 2\frac{h\nu^3}{c^2}\frac{1}{e^\frac{h\nu}{\kb T}-1}\vec{n}\cdot\vec{A}d\Omega d\nu
\end{eqnarray*}
$P$ ist die Leistung, die bei der Frequenz $\nu$ durch ein Flächenelement $\vec{A}$ und einen Raumwinkel $\Omega$ in Richtung $\vec{n}$ bei der Temperatur $T$ von einem Schwarzen Strahler abgestrahlt wird.
\begin{eqnarray*}
I(\nu,T) &=& 2\frac{h\nu^3}{c^2}\frac{1}{e^\frac{h\nu}{\kb T}-1}
\end{eqnarray*}
heißt spektrale Intensität. Über den Halbraum integriert:
\begin{eqnarray*}
\int I(\nu,T) \vec{n}\cdot\vec{A}d\Omega d\nu
&=& I(\nu,T)d\nu \int \vec{n}\cdot\vec{A}d\Omega \\
&=& I(\nu,T)d\nu A\int_0^\frac{\pi}{2} \cos(\phi)2\pi\sin(\phi) d\phi \\
&=& I(\nu,T)d\nu A\pi\int_0^\frac{\pi}{2} \sin(2\phi) d\phi \\
&=& I(\nu,T)d\nu A\pi\left(-\frac{1}{2}\cos(2\phi)\right|_0^\frac{\pi}{2}\\
&=& I(\nu,T)d\nu A\pi
\end{eqnarray*}


\end{subsection}

\begin{subsection}{Gesamt-Energie-Dichte}
Integriert man über die Frequenz $\nu$ erhält man die Gesamt-Energie:
\begin{eqnarray*}
E(V,T) &=& \frac{8h\pi}{3}\frac{V}{c^3} \int_0^\infty\nu^3 \frac{1}{e^{\beta h\nu}-1}d\nu
\end{eqnarray*}
Wir ersetzen $x=\beta h\nu$, $\nu=\frac{x}{\beta h}=\frac{\kb T}{h}x$
\begin{eqnarray*}
E(V,T) &=& \frac{8h\pi}{3}\frac{V}{c^3} \left(\frac{\kb T}{h}\right)^4\int_0^\infty  \frac{x^3}{e^x-1}dx
\end{eqnarray*}
Dabei ist
\begin{eqnarray*}
\int_0^\infty  \frac{x^3}{e^x-1}dx &=& \frac{\pi^4}{15}
\end{eqnarray*}
\begin{eqnarray*}
E(V,T) &=& \frac{8h\pi}{3}\frac{V}{c^3} \left(\frac{\kb T}{h}\right)^4\frac{\pi^4}{15}\\
&=& \frac{8\pi^5\kb^4}{h^3c^3 3\cdot 15}T^4V
\end{eqnarray*}
Analog ist die Gesamt-Strahlungsdichte
\begin{eqnarray*}
\int_0^\infty I(\nu,T)d\nu &=& \int_0^\infty 2\frac{h\nu^3}{c^2}\frac{1}{e^\frac{h\nu}{\kb T}-1}d\nu\\
&=& 2\frac{h}{c^2}\left(\frac{\kb T}{h}\right)^4\int_0^\infty \frac{x^3}{e^x-1}dx\\
&=& \frac{2(T\kb)^4}{h^3c^2}\frac{\pi^4}{15}
\end{eqnarray*}

Dies ist das Stefan-Boltzmann-Gesetz.
\end{subsection}

\begin{subsection}{Wiensches Verschiebungsgesetz}
Wir betrachten die Energiedichte in der Wellenlänge ausgedrückt \ref{lambda_density}:
\begin{eqnarray*}
d E(\lambda,V,T) &=&\frac{8h\pi}{3}\frac{Vc}{\lambda^5} \frac{1}{e^{\frac{c\beta h}{\lambda}}-1}d\lambda
\end{eqnarray*}
Um die maximale Energiedichte zu ermitteln können Vorfaktoren ignorierte werden. Wir setzen
\begin{eqnarray*}
 x &=& \frac{\lambda}{c\beta h}
\end{eqnarray*}
und leiten ab:
\begin{eqnarray*}
\frac{d}{dx}\frac{1}{x^5\left(e^{x^{-1}}-1\right)}&=&0
\end{eqnarray*}
Dies hat aber dieselbe Lösung wie
\begin{eqnarray*}
\frac{d}{dx}\left(x^5\left(e^{x^{-1}}-1\right)\right)&=&0\\
&=& 5x^4\left(e^{x^{-1}}-1\right)-x^5x^{-1}x^{-2}e^{x^{-1}}\\
&=&\left(5x^4-x^2\right)e^{x^{-1}}-5x^3\\
1&=&\left(x-\frac{1}{5x}\right)e^{x^{-1}}
\end{eqnarray*}
\end{subsection}

\begin{subsection}{Nullpunkt-Energie}
Eine exakte quantenmechanische Behandlung des harmonischen Oszillators führt zu Energie-Niveaus 
\begin{eqnarray*}
E_n(\nu) &=& h\nu \left(\frac{1}{2}+n\right)
\end{eqnarray*}
Die führt erneut zu eine Ultraviolett-Katastrophe!!
\end{subsection}
\end{section}

\begin{section}{Bohrs Modell des Wasserstoff-Atoms}
\begin{subsection}{Fraunhoferlinien}
\end{subsection}
\begin{subsection}{Elektronen auf Planeten-Bahnen}
Die Kraft zwischen zwei Ladungen ist
\begin{eqnarray*}
F &=& \frac{1}{4\pi\epsilon_0}\frac{q_1q_2}{r^2}
\end{eqnarray*}
Für das Wasserstoff-Atom:
\begin{eqnarray*}
F &=& \frac{1}{4\pi\epsilon_0}\frac{e^2}{r^2}
\end{eqnarray*}
Dies muss gleich der Zentripetal-Beschleunigung mal die Elektronen-Masse $m$ sein.
\begin{eqnarray*}
F &=& mr\omega^2
\end{eqnarray*}
Dies führt zu einem Zusammenhang zwischen $r$ und $\omega$:
\begin{eqnarray*}
mr\omega^2&=& \frac{1}{4\pi\epsilon_0}\frac{e^2}{r^2}\\
\omega^2&=& \frac{1}{4\pi\epsilon_0}\frac{e^2}{mr^3}
\end{eqnarray*}
\end{subsection}

\begin{subsection}{Quantisierungs-Bedingung}
Der Drehimpuls $L$ hat sie selbe physikalische Einheit wie das Plancksche Wirkungs-Quantum $\hbar$.
Für das Wasserstoff-Atom setzte Bohr an:
\begin{eqnarray}
\label{bohr_condition}
L &=& n\hbar = mr^2\omega
\end{eqnarray}
oder
\begin{eqnarray*}
\omega &=& \frac{n\hbar}{mr^2}
\end{eqnarray*}
Das bestimmt den Radius $r$:
\begin{eqnarray*}
\frac{1}{4\pi\epsilon_0}\frac{e^2}{mr^3} &=& \left(\frac{n\hbar}{mr^2}\right)^2\\
\frac{1}{4\pi\epsilon_0}\frac{e^2}{mr^3} &=& \left(\frac{n\hbar}{m}\right)^2\frac{1}{r^4}\\
\frac{1}{4\pi\epsilon_0}\frac{e^2r}{m} &=& \left(\frac{n\hbar}{m}\right)^2\\
r &=& \left(\frac{n\hbar}{m}\right)^2\frac{4\pi\epsilon_0m}{e^2}\\
&=&n^2\frac{4\pi\hbar^2\epsilon_0}{e^2m}
\end{eqnarray*}
$a_0r$ heißt bohrsche Radius.
\end{subsection}

\begin{subsection}{Energie-Niveaus}
Die potentielle Energie ist
\begin{eqnarray*}
V_n &=& -\frac{e^2}{4\pi\epsilon_0}\frac{1}{r}=-\frac{e^2}{4\pi\epsilon_0}\frac{1}{n^2}\frac{e^2m}{4\pi\hbar^2\epsilon_0}\\
&=& -\frac{1}{n^2}\frac{e^4m}{(4\pi\epsilon_0\hbar)^2}
\end{eqnarray*}
Für die kinetische Energie erhalten wir
\begin{eqnarray*}
T_n &=& \frac{1}{2}mv^2 =\frac{1}{2}m(r\omega)^2\\
(r\omega)^2 &=& r^2\frac{1}{4\pi\epsilon_0}\frac{e^2}{mr^3}\\
&=& \frac{1}{4\pi\epsilon_0}\frac{e^2}{mr}\\
&=& \frac{1}{4\pi\epsilon_0}\frac{e^2}{m}\frac{1}{n^2}\frac{e^2m}{4\pi\hbar^2\epsilon_0}\\
&=&\frac{1}{n^2}\frac{e^4}{(4\pi\epsilon_0\hbar)^2}\\
T_n &=& \frac{1}{2}\frac{1}{n^2}\frac{e^4m}{(4\pi\epsilon_0\hbar)^2}
\end{eqnarray*}
Und damit die Gesamt-Energie
\begin{eqnarray*}
E_n &=& -\frac{1}{2}\frac{1}{n^2}\frac{e^4m}{(4\pi\epsilon_0\hbar)^2} = -\frac{1}{2}\frac{1}{n^2}\frac{e^4m}{(2\epsilon_0 h)^2} 
\end{eqnarray*}
$E_1$ heißt auch Rydberg-Energie $Ry$:
\begin{eqnarray*}
E_n &=& -\frac{1}{n^2}Ry
\end{eqnarray*}

\end{subsection}
\end{section}

\begin{section}{Louis de Broglie}
\begin{subsection}{Wellen-Teilchen-Dualismus}
\end{subsection}
\begin{subsection}{de Broglie-Hypothese}
De Broglie nahm an, dass jeder Energie eine Frequenz zukommt:
\begin{eqnarray*}
E &=& h\nu = \hbar\omega\\
\omega &=& \frac{E}{\hbar}
\end{eqnarray*}
Also auch einem Teilchen der Ruhemasse $m$
\begin{eqnarray*}
\omega_0(m) &=& \frac{mc^2}{\hbar}
\end{eqnarray*}
\textit{"Die Kühnheit dieser Idee war so groß – ich muss aufrichtig sagen, dass ich selber auch damals den Kopf schüttelte dazu, und ich erinnere mich sehr gut, dass Herr Lorentz mir damals sagte im vertraulichen Privatgespräch: 'Diese jungen Leute nehmen es doch gar zu leicht, alte physikalische Begriffe beiseite zu setzen!' Es war damals die Rede von Broglie-Wellen, von der heisenbergschen Unschärfe-Relation – das schien damals uns Älteren etwas sehr schwer Verständliches."'}
\end{subsection}

\begin{subsection}{Wellen in relativistischer Raum-Zeit}
Eine Welle in einer Dimension hat typischerweise sie Form
\begin{eqnarray*}
\phi(t,x) &=& \sin(\omega t-kx)
\end{eqnarray*}
wobei $k$ die Ortsfrequenz auch Wellenzahl genannt ist. Die Ausbreitungsgeschwindigkeit $v$ ist
\begin{eqnarray*}
v &=& \frac{\omega}{k}
\end{eqnarray*}
denn
\begin{eqnarray*}
\omega(t+\Delta t)-k(x+v\Delta t) &=& \omega t-kx
\end{eqnarray*}
In relativistischen Ereignis-Koordinaten:
\begin{eqnarray*}
\phi(x^\mu) &=& \sin\left(\frac{\omega}{c}x^0-\vec{k}\cdot\vec{x}\right)
\end{eqnarray*}
oder in einer Raum-Zei mit nur einer Raum-Koordinate
\begin{eqnarray*}
\phi(x^\mu) &=& \sin\left(\frac{\omega}{c}x^0-kx\right)\\
k_\mu &=& \pmat{\frac{\omega}{c} & -k}\\
\phi(x^\mu) &=& \sin\left(k_\mu x^\mu\right)
\end{eqnarray*}
$k_\mu$ transformiert sich demnach wie ein dualer Vektor.
\end{subsection}

\begin{subsection}{Materie-Wellen}
Ein ruhendes Teilchen der Masse $m$ entspricht eine stehenden Welle
\begin{eqnarray*}
\phi(x^\mu) &=& \sin\left(\frac{\omega_0(m)}{c}x^0\right)\\
 &=& \sin\left(\frac{mc}{\hbar}x^0\right)
\end{eqnarray*}
Dies entspricht einem Wellenvektor
\begin{eqnarray*}
k_\mu &=& \pmat{\frac{mc}{\hbar} & 0}
\end{eqnarray*}
Duale Vektoren transformieren sich mit der inversen Matrix:
\begin{eqnarray*}
k^\prime_\mu &=& \pmat{\frac{mc}{\hbar} & 0}\gamma(v)
\pmat{1 & -\frac{v}{c}\\
 -\frac{v}{c} &1
}\\
&=& \gamma(v)\pmat{\frac{mc}{\hbar} & -\frac{mv}{\hbar}}
\end{eqnarray*}
Die Ausbreitungsgeschwindigkeit ist
\begin{eqnarray*}
v &=& \frac{\omega}{k} = \frac{c}{v}
\end{eqnarray*}
größer als die Lichtgeschwindigkeit! Die Wellenlänge $\lambda$ ergibt sich
\begin{eqnarray}
\label{de_broglie_wavelenght}
k\lambda &=& 2\pi\cr
\lambda &=& \frac{2\pi}{k}\cr
\lambda(m,v) &=& \frac{2\pi}{\gamma(v)\frac{mv}{\hbar}}=\frac{2\pi\hbar}{\gamma(v)mv}\cr
&=&\frac{h}{p}
\end{eqnarray}
Denn $p=\gamma(v)mv$ ist der relativistische Ausdruck für den Impuls eines Teilchens.
$\lambda=\frac{h}{p}$ heißt die de Broglie Wellenlänge.
\\ \\
Wikipedia:

\textit{Einstein zeigte sich tief beeindruckt und erklärte später, er glaube, dass de Broglies Doktorarbeit den ersten schwachen Lichtstrahl auf dieses leidigste unter den physikalischen Rätseln werfe.
}
\end{subsection}

\begin{subsection}{Bohrs Quantisierungs-Bedingung}
Bohr Quantisierungs-Bedingung \ref{bohr_condition} war
\begin{eqnarray*}
L &=& n\hbar = mr^2\omega
\end{eqnarray*}
De Broglie hingegen argumentierte anders: Der Bahnumfang muss ein Vielfaches der Wellenlänge sein:
\begin{eqnarray*}
2\pi r &=& n\lambda\\
\lambda=\frac{h}{p}\\
p &=& mr\omega\\
2\pi r &=& n\frac{h}{rm\omega}\\
r^2m\omega &=& n\hbar
\end{eqnarray*}
Die beiden Bedingungen sind Äquivalent.
\end{subsection}

\begin{subsection}{Davisson-Germer-Experiment}
\end{subsection}
\end{section}

\begin{section}{Schrödinger-Gleichung}
\begin{subsection}{Zeitunabhängige Variante}
Erwin Schrödinger hat sich die Argumente de Broglie's zu Herzen genommen. Dabei war zunächst die de Broglie-Wellenlänge \ref{de_broglie_wavelenght} der Ausgangspunkt.
\begin{eqnarray*}
\lambda &=&\frac{h}{p}
\end{eqnarray*}
Damit hängt auch die Ortsfrequenz $k$ zusammen:
\begin{eqnarray*}
k &=& \frac{2\pi}{\lambda} = \frac{2\pi p}{h}=\frac{p}{\hbar}
\end{eqnarray*}
Dies kann auch als Vektor-Gleichung verstanden werden
\begin{eqnarray*}
\vec{k} &=&  \frac{\vec{p}}{\hbar}
\end{eqnarray*}
Für eine ebene Welle bedeutet dies:
\begin{eqnarray*}
\phi(t,\vec{x}) &=& \sin(\omega t-\vec{k}\cdot\vec{p})\\
\vec{\nabla}\phi(t,\vec{x}) &=& -\vec{k}\cos(\omega t-\vec{k}\cdot\vec{p})\\
\Delta\phi(t,\vec{x}) &=& -\vec{k}^2\sin(\omega t-\vec{k}\cdot\vec{p})=-\left(\frac{\vec{p}}{\hbar}\right)^2\sin(\omega t-\vec{k}\cdot\vec{p})
\end{eqnarray*}
Die Zeit $t$ spielt hier keine Rolle, auf eine beliebige Funktion $\phi(\vec{x})$ angewandt:
\begin{eqnarray*}
\Delta \phi(\vec{x}) &=&-\left(\frac{\vec{p}}{\hbar}\right)^2\phi(\vec{x})\\
\vec{p}^2\phi(\vec{x}) &=&-\hbar^2\Delta \phi(\vec{x})
\end{eqnarray*}
Daher kann die kinetische Energie im Wasserstoff-Atom so geschrieben werden:
\begin{eqnarray*}
T\phi(\vec{x})&=&\frac{\vec{p}^2}{2m}\phi(\vec{x})=-\frac{\hbar^2}{2m}\Delta\phi(\vec{x})
\end{eqnarray*}
Nehmen wir nun auch noch die Potentielle Energie hinzu erhalten wir
\begin{eqnarray*}
E\phi(\vec{x})&=&\left(-\frac{\hbar^2}{2m}\Delta-\frac{e^2}{4\pi\epsilon_0}\frac{1}{|\vec{x}|}\right)\phi(\vec{x})
\end{eqnarray*}
Das ist die sogenannte zeitunabhängige Schrödinger-Gleichung für ein Wasserstoff-Atom.
Der in Klammern stehende Ausdruck wird Hamilton-Operator genannt:
\begin{eqnarray}
\label{hydrogen_hamilton_operator_def}
\hat{H}&=&-\frac{\hbar^2}{2m}\Delta-\frac{e^2}{4\pi\epsilon_0}\frac{1}{|\vec{x}|}
\end{eqnarray}
Also
\begin{eqnarray*}
E \phi(\vec{x}) &=&\hat{H}\phi(\vec{x})
\end{eqnarray*}
Sie stellt ein Eigenwertproblem dar. Dieses kann sogar analytisch gelöst werden:
\begin{eqnarray*}
\phi_{nlm}(\vec{x})&=& R_{nl}(|\vec{r}|)Y_{lm}\left(\frac{\vec{x}}{|\vec{x}|}\right)
\end{eqnarray*}
$n$, $l$ und $m$ sind drei ganzzahlige Quantenzahlen. Und zwar so:
\begin{eqnarray*}
n &>& 0\\
l &=& 0..(n-1)\\
m&=& -l .. l
\end{eqnarray*}
Die Energie $E$ ist gleich den Eigenwerten und nehmen dieselben Werte an, wie beim Bohr-Model:
\begin{eqnarray*}
-\frac{1}{n^2}Ry \phi_{nlm}(\vec{x}) &=&\hat{H}\phi_{nlm}(\vec{x})
\end{eqnarray*}
Die $Y_{nlm}$ sind komplexe Funktionen.
\begin{eqnarray*}
Y_{lm}\left(\pmat{ \cos(\varphi)\sin(\theta)
\\
\sin(\varphi)\sin(\theta)\\
\cos(\theta) }\right) &=& p(\cos(\theta)e^{im\varphi}
\end{eqnarray*}
können aber auch durch reelle Funktionen ersetzt werden:
\begin{eqnarray*}
Y^c_{lm}\left(\pmat{ \cos(\varphi)\sin(\theta)
\\
\sin(\varphi)\sin(\theta)\\
\cos(\theta) }\right) &=& p(\cos(\theta)\cos(m\varphi)\\
Y^s_{lm}\left(\pmat{ \cos(\varphi)\sin(\theta)
\\
\sin(\varphi)\sin(\theta)\\
\cos(\theta) }\right) &=& p(\cos(\theta)\sin(m\varphi)\\
Y_{nlm} &=& Y^c_{nlm}+iY^s_{nlm}
\end{eqnarray*}
Die $Y^c_{nlm}$ und $Y^s_{nlm}$ bilden ein vollständiges Funktionen-System. Das heißt, jede beliebige reelle Funktion kann als Linearkombination dieser Funktionen dargestellt werden.
\end{subsection}

\begin{subsection}{Zeitabhängige, reelle Variante}
Mit der zeitunabhängigen Variante wird der de Broglie-Beziehung
\begin{eqnarray*}
\lambda &=&\frac{h}{p}
\end{eqnarray*}
Rechnung getragen. Es wäre nun schön, wenn auch die Planck-Beziehung
\begin{eqnarray*}
\nu &=& \frac{E}{h}\\
\omega &=& \frac{E}{\hbar}
\end{eqnarray*}
in die Gleichung einbezogen werden könnte. Noch schöner wäre es, wenn wir einen Operator finden könnten der die Dynamik der Wellenfunktion beschreibt:
\begin{eqnarray*}
\frac{\partial}{\partial t} \phi(t,\vec{x}) &=& \hat{O}\phi(t,\vec{x})
\end{eqnarray*}
Dabei stellt sich aber ein Problem. Nehmen wir an für eine Funktion $\phi(\vec{x})$ gelte
\begin{eqnarray*}
E\phi(\vec{x}) &=& \hat{H}\phi(\vec{x})
\end{eqnarray*}
Dann könnten wir versucht sein die zeitabhängige Welle
\begin{eqnarray*}
\phi(t,\vec{x}) &=& \sin(\omega t)\phi(\vec{x})=\sin\left(\frac{E}{\hbar} t\right)\phi(\vec{x})
\end{eqnarray*}
in Betracht zu ziehen. Nach der Zeit abgeleitet:
\begin{eqnarray*}
\frac{\partial}{\partial t} \phi(t,\vec{x}) &=& \omega\cos(\omega t)\phi(\vec{x})
\end{eqnarray*}
ähnlich verhält es sich mit
\begin{eqnarray*}
\phi(t,\vec{x}) &=& \cos(\omega t)\phi(\vec{x})\\
\frac{\partial}{\partial t} \phi(t,\vec{x}) &=& -\omega\sin(\omega t)\phi(\vec{x})
\end{eqnarray*}
Eine Lösung besteht darin, anstelle einer einzigen Funktion, ein Paar von Funktionen zu betrachten zum Beispiel:
\begin{eqnarray*}
\phi(t,\vec{x}) &=& \pmat{\cos(\omega t)\\ \sin(\omega t)}\phi(\vec{x})\\
\frac{\partial}{\partial t} \phi(t,\vec{x}) &=& \omega\pmat{-\sin(\omega t)\\ \cos(\omega t)}\phi(\vec{x})
\end{eqnarray*}
Nun müssen nur noch $\sin$ und $\cos$ miteinander vertauscht werden. Dies kann mit Hilfe einer $2\times 2$-Matrix erreicht werden:
\begin{eqnarray*}
\frac{\partial}{\partial t} \pmat{ 0 & 1\\ -1 & 0}\phi(t,\vec{x}) &=& \omega\pmat{ 0 & 1\\ -1 & 0}\pmat{-\sin(\omega t)\\ \cos(\omega t)}\phi(\vec{x})\\
&=& \omega\pmat{\cos(\omega t)\\ \sin(\omega t)}\phi(\vec{x})\\
&=& \omega \phi(t,\vec{x}) =\frac{E}{\hbar} \phi(t,\vec{x}) =\frac{\hat{H}}{\hbar} \phi(t,\vec{x}) 
\end{eqnarray*}
Nun gilt
\begin{eqnarray*}
\pmat{ 0 & 1\\ -1 & 0}\pmat{ 0 & 1\\ -1 & 0} &=& -\pmat{ 1 & 0\\ 0 & 1}
\end{eqnarray*}
Daher kann obige Gleichung umgeschrieben werden:
\begin{eqnarray}
\label{schroedinger_time_1}
\frac{\partial}{\partial t} \phi(t,\vec{x}) &=&=-\pmat{ 0 & 1\\ -1 & 0}\frac{\hat{H}}{\hbar} \phi(t,\vec{x}) 
\end{eqnarray}
und wir erhalten die gewünschte Dynamik der Wellenfunktion. Diese soll nun auf eine beliebige Funktion
\begin{eqnarray*}
\phi(t,\vec{x}) &=&= \pmat{ \phi_1(t, \vec{x})\\ \phi_2(t, \vec{x})}
\end{eqnarray*}
angewandt werden können. Nun kann aber jedes beliebige $\phi(t, \vec{x})$ so dargestellt werden:
\begin{eqnarray*}
\phi(t, \vec{x})=\pmat{\phi_1(t, \vec{x})\\ \phi_2(t, \vec{x})} &=& 
\pmat{
\sum_{nlm} \tau^c_{1,nlm}(t)\phi^c_{nlm}(\vec{x})+\tau^s_{1,nlm}(t)\phi^s_{nlm}(\vec{x})\\
\sum_{nlm} \tau^c_{2,nlm}(t)\phi^c_{nlm}(\vec{x})+\tau^s_{2,nlm}(t)\phi^s_{nlm}(\vec{x})
}
\end{eqnarray*}
Wenden wir darauf den Hamilton-Operator an:
\begin{eqnarray*}
\hat{H}\phi(t, \vec{x})=\pmat{\phi_1(t, \vec{x})\\ \phi_2(t, \vec{x})} &=& 
\pmat{
\sum_{n}E_n\sum_{lm} \tau^c_{1,nlm}(t)\phi^c_{nlm}(\vec{x})+\tau^s_{1,nlm}(t)\phi^s_{nlm}(\vec{x})\\
\sum_{n}E_n\sum_{lm} \tau^c_{2,nlm}(t)\phi^c_{nlm}(\vec{x})+\tau^s_{2,nlm}(t)\phi^s_{nlm}(\vec{x})
}
\end{eqnarray*}
mit
\begin{eqnarray*}
E_n &=& -\frac{1}{n^2}Ry
\end{eqnarray*}
und weiter
\begin{eqnarray*}
-\pmat{ 0 & 1\\ -1 & 0}
\frac{\hat{H}}{\hbar}\phi(t, \vec{x})&=& \frac{1}{\hbar}
\pmat{
-\sum_{n}E_n\sum_{lm} \tau^c_{2,nlm}(t)\phi^c_{nlm}(\vec{x})+\tau^s_{2,nlm}(t)\phi^s_{nlm}(\vec{x})\\
\sum_{n}E_n\sum_{lm} \tau^c_{1,nlm}(t)\phi^c_{nlm}(\vec{x})+\tau^s_{1,nlm}(t)\phi^s_{nlm}(\vec{x})
}
\end{eqnarray*}
dann haben wir die rechte Seite von \ref{schroedinger_time_1}. Die linke Seite ergibt:
\begin{eqnarray*}
\frac{\partial}{\partial t}\phi(t, \vec{x})=\pmat{\phi_1(t, \vec{x})\\ \phi_2(t, \vec{x})} &=& 
\pmat{
\sum_{nlm} \tau^{c\prime}_{1,nlm}(t)\phi^c_{nlm}(\vec{x})+\tau^{s\prime}_{1,nlm}(t)\phi^s_{nlm}(\vec{x})\\
\sum_{nlm} \tau^{c\prime}_{2,nlm}(t)\phi^c_{nlm}(\vec{x})+\tau^{s\prime}_{2,nlm}(t)\phi^s_{nlm}(\vec{x})
}
\end{eqnarray*}
Setzen wir die beiden Seiten gleich, dann müssen die Koeffizienten übereinstimmen, denn die $\phi^c_{nlm}$ und $\phi^s_{nlm}$ sind alle voneinander linear unabhängig:
\begin{eqnarray*}
-\frac{E_n}{\hbar}\tau^c_{2,nlm}(t)&=&\tau^{c\prime}_{1,nlm}(t)\\
\frac{E_n}{\hbar}\tau^c_{1,nlm}(t)&=&\tau^{c\prime}_{2,nlm}(t)\\
-\frac{E_n}{\hbar}\tau^s_{2,nlm}(t)&=&\tau^{s\prime}_{1,nlm}(t)\\
\frac{E_n}{\hbar}\tau^s_{1,nlm}(t)&=&\tau^{s\prime}_{2,nlm}(t)
\end{eqnarray*}
mit
\begin{eqnarray*}
\omega_n &=& \frac{E_n}{\hbar}
\end{eqnarray*}
also
\begin{eqnarray*}
-\omega_n^2\tau^c_{2,nlm}(t)&=&\tau^{c\prime\prime}_{2,nlm}(t)\\
-\omega_n\tau^s_{2,nlm}(t)&=&\tau^{s\prime\prime}_{2,nlm}(t)
\end{eqnarray*}
Ganz allgemein ist die Lösung von
\begin{eqnarray*}
-\omega^2\tau(t) = \tau^{\prime\prime}(t)&\Leftrightarrow& \tau(t) = a\sin(\omega t+\xi)
\end{eqnarray*}
also
\begin{eqnarray*}
\tau^c_{2,nlm}(t)&=&a^c_{nlm}\sin(\omega_n t+\xi^c_{nlm})\\
\tau^s_{2,nlm}(t)&=&a^s_{nlm}\sin(\omega_n t+\xi^s_{nlm})
\end{eqnarray*}
und daher
\begin{eqnarray*}
\tau^c_{1,nlm}(t)&=&a^c_{nlm}\cos(\omega_n t+\xi^c_{nlm})\\
\tau^s_{1,nlm}(t)&=&a^s_{nlm}\cos(\omega_n t+\xi^s_{nlm})
\end{eqnarray*}
Das heißt, die allgemeine Lösung von \ref{schroedinger_time_1} lautet
\begin{eqnarray*}
\phi(t, \vec{x})=\pmat{\phi_1(t, \vec{x})\\ \phi_2(t, \vec{x})} &=& \sum_{nlm} \left(a^c_{nlm}\pmat{\cos(\omega_n t+\xi^c_{nlm})\\\sin(\omega_n t+\xi^c_{nlm})}\phi^c_{nlm}
+a^s_{nlm}\pmat{\cos(\omega_n t+\xi^s_{nlm})\\\sin(\omega_n t+\xi^s_{nlm})}\phi^s_{nlm}\right)
\end{eqnarray*}
In anderen Worten, die Anteile zu unterschiedlichen Energie-Eigenwerten $E_n$ schwingen unabhängig voneinander mit der entsprechenden Frequenz $\omega_n$.
\end{subsection}

\begin{subsection}{Zeitabhängige, komplexe Variante}
Dasselbe wie im letzten Abschnitt lässt sich mit komplexen Zahlen beschreiben. Die Matrix $\pmat{ 0 & 1\\ -1 & 0}$ hat, wie oben schon erwähnt, die Eigenschaft
\begin{eqnarray*}
\pmat{ 0 & 1\\ -1 & 0}\pmat{ 0 & 1\\ -1 & 0} &=& -\pmat{ 1 & 0\\ 0 & 1}
\end{eqnarray*}
Komplexe Zahlen können als Matrizen
\begin{eqnarray*}
a+bi &=& \pmat{ a & b\\ -b & a}
\end{eqnarray*}
aufgefasst werden. Das heißt
\begin{eqnarray*}
i &=& \pmat{ 0 & 1\\ -1 & 0}
\end{eqnarray*}
Das heißt
\begin{eqnarray*}
i^2 &=& -1
\end{eqnarray*}
Statt
\begin{eqnarray*}
\phi(t,\vec{x})&=&\pmat{\phi_1(t,\vec{x})\\ \phi_2(t,\vec{x})}
\end{eqnarray*}
schreiben wir
\begin{eqnarray*}
\phi(t,\vec{x})&=&\phi_2(t,\vec{x})+i\phi_1(t,\vec{x})
\end{eqnarray*}
denn
\begin{eqnarray*}
\pmat{ 0 & 1\\ -1 & 0}\pmat{\phi_1(t,\vec{x})\\ \phi_2(t,\vec{x})}
&=&
\pmat{\phi_2(t,\vec{x})\\-\phi_1(t,\vec{x})}\\
i(\phi_2(t,\vec{x})+i\phi_1(t,\vec{x})) &=& -\phi_1(t,\vec{x})+i\phi_2(t,\vec{x})
\end{eqnarray*}
\ref{schroedinger_time_1} wird dann zu
\begin{eqnarray}
\label{schrödinger_complex}
\frac{\partial}{\partial t} \phi(t,\vec{x}) &=&-i\frac{\hat{H}}{\hbar} \phi(t,\vec{x}) \cr
i\hbar\frac{\partial}{\partial t} \phi(t,\vec{x}) &=&\hat{H}\phi(t,\vec{x})
\end{eqnarray}
In dieser Form erscheint die zeitabhängige Schrödinger-Gleichung in der Literatur.

Die allgemeine Welle, die \ref{schrödinger_complex} erfüllt, ist
\begin{eqnarray*}
\phi(t,\vec{x}) &=& \sum_{n}e^{-i\omega_n t}\sum_{lm} a_{nlm}Y_{nlm}(\vec{x})
\end{eqnarray*}
wobei $a_{nlm}$ komplexe Zahlen sind, denn
\begin{eqnarray*}
i\hbar\frac{\partial}{\partial t}\phi(t,\vec{x}) &=& i\hbar\sum_{n}-i\omega_ne^{-i\omega_n t}\sum_{lm} a_{nlm}\phi_{nlm}(\vec{x})\\
&=& \hbar\sum_{n}\omega_ne^{-i\omega_n t}\sum_{lm} a_{nlm}\phi_{nlm}(\vec{x})\\
&=& \sum_{n}E_ne^{-i\omega_n t}\sum_{lm} a_{nlm}\phi_{nlm}(\vec{x})\\
&=& \hat{H}\phi(t,\vec{x})
\end{eqnarray*}
Für ein bestimmtes Energie-Niveau $E_n$
\begin{eqnarray*}
\phi_n(t,\vec{x}) &=& e^{-i\omega_n t}\sum_{lm} a_{nlm}\phi_{nlm}(\vec{x})
\end{eqnarray*}
\end{subsection}

\begin{subsection}{Interpretation}
\begin{subsubsection}{Vorbereitung}
Zunächst führen wir eine Notation ein:
\begin{eqnarray*}
\braket{\phi(t,\vec{x})}{\xi(t,\vec{x})}&=&\int dx^3 \phi^*(t,\vec{x})\xi(t,\vec{x})
\end{eqnarray*}
Für normierte Wellenfunktionen gilt
\begin{eqnarray*}
\braket{\phi(t,\vec{x})}{\phi(t,\vec{x})}&=&1
\end{eqnarray*}
Wir betrachten zwei Raum-Funktionen $\Phi_1(\vec{x})$ und $\Phi_2(\vec{x})$ mit den Eigenschaften
\begin{eqnarray*}
\hat{H}\Phi_1(\vec{x}) &=& E_1\Phi_1(\vec{x})\\
\braket{\Phi_1(\vec{x})}{\Phi_1(\vec{x})} &=& 1\\
\hat{H}\Phi_2(\vec{x}) &=& E_2\Phi_2(\vec{x})\\
\braket{\Phi_2(\vec{x})}{\Phi_2(\vec{x})} &=& 1
\end{eqnarray*}
Wobei $E_1\ne E_2$. Daraus lassen sich zwei Lösungen der zeitabhängigen Schrödinger-Gleichung \ref{schrödinger_complex}
\begin{eqnarray*}
i\hbar\frac{\partial}{\partial t} \phi(t,\vec{x}) &=&\hat{H}\phi(t,\vec{x})
\end{eqnarray*}
bilden:
\begin{eqnarray*}
\phi_1(t,\vec{x}) &=& e^{-i\omega_1 t}\Phi_1(\vec{x})\\
\phi_2(t,\vec{x}) &=& e^{-i\omega_2 t}\Phi_2(\vec{x})
\end{eqnarray*}
mit
\begin{eqnarray*}
\omega_1 &=& \frac{E_1}{\hbar}\\
\omega_2 &=& \frac{E_2}{\hbar}
\end{eqnarray*}
Nun bilden wir eine Linearkombination $\phi$ von $\phi_1$ und $\phi_2$:
\begin{eqnarray}
\label{double_wave_function}
\phi(t,\vec{x}) &=& a_1\phi_1(t,\vec{x})+a_2\phi_2(t,\vec{x})
\end{eqnarray}
wobei für die beiden komplexen Zahlen $a_1$ und $a_2$
\begin{eqnarray*}
a_1^*a_1+a_2^*a_2 &=& 1
\end{eqnarray*}
gilt. Das hat zur Folge, dass auch $\phi$ normiert ist:
\begin{eqnarray*}
\braket{\phi(t,\vec{x})}{\phi(t,\vec{x})} &=& 1
\end{eqnarray*}
Denn
\begin{eqnarray*}
\braket{\phi(t,\vec{x})}{\phi(t,\vec{x})} &=& \braket{a_1\phi_1(t,\vec{x})+a_2\phi_2(t,\vec{x})}{a_1\phi_1(t,\vec{x})+a_2\phi_2(t,\vec{x})}\\
&=&a_1^*a_1\braket{\phi_1(t,\vec{x})}{\phi_1(t,\vec{x})}\\
&&+a_1^*a_2\braket{\phi_1(t,\vec{x})}{\phi_2(t,\vec{x})}\\
&&+a_2^*a_1\braket{\phi_2(t,\vec{x})}{\phi_1(t,\vec{x})}\\
&&+a_2^*a_2\braket{\phi_2(t,\vec{x})}{\phi_2(t,\vec{x})}\\
&=&a_1^*a_1\braket{\Phi_1(t,\vec{x})}{\Phi_1(t,\vec{x})}\\
&&+a_2^*a_2\braket{\Phi_2(t,\vec{x})}{\Phi_2(t,\vec{x})}\\
&=&a_1^*a_11+a_2^*a_21 = 1
\end{eqnarray*}
$\braket{\phi_1(t,\vec{x})}{\phi_2(t,\vec{x})}=0$ und $\braket{\phi_2(t,\vec{x})}{\phi_1(t,\vec{x})}=0$ (siehe \ref{eigen_normality}).
Daher kann, wie oben schon angesprochen, $\phi^*(t,\vec{x})\phi(t,\vec{x})$ als Aufenthaltswahrscheinlichkeit und
\begin{eqnarray*}
\rho(t,\vec{x}) &=& -e\phi^*(t,\vec{x})\phi(t,\vec{x})
\end{eqnarray*}
als Ladungsdichte aufgefasst werden. Diese ist zwar immer noch reell und positiv, aber nicht mehr konstant.
\end{subsubsection}

\begin{subsubsection}{\titlemath{$\hat{H}$}{H} ist ein hermitescher Operator}
Für einen beliebigen Operator $\hat{A}$ ist der dazu adjungierte Operator $\hat{A}^\dagger$ durch
\begin{eqnarray*}
\braket{\hat{A}^\dagger \phi}{\chi} &=& \braket{\phi}{\hat{A}\chi}
\end{eqnarray*}
(gültig für beliebige $\phi$ und $\chi$) definiert. Gilt $\hat{A}^\dagger = \hat{A}$, spricht man von einem selbstadjungierten oder hermiteschen Operator. Eigenwerte hermitescher Operatoren sind reell:
\begin{eqnarray*}
\hat{A}\phi(\vec{x}) &=& a\phi(\vec{x})\\
\braket{\hat{A}\phi}{\phi} &=& \braket{\phi}{\hat{A}\phi}\\
\braket{a\phi}{\phi} &=& \braket{\phi}{a\phi}\\
a^*\braket{\phi}{\phi} &=& a\braket{\phi}{\phi}\\
a^* &=& a
\end{eqnarray*}


Wir wollen nun zeigen, dass der Hamilton-Operator \ref{hydrogen_hamilton_operator_def} in der Schrödinger-Gleichung für das Wasserstoff-Atom.
\begin{eqnarray*}
\hat{H}&=&-\frac{\hbar^2}{2m}\Delta-\frac{e^2}{4\pi\epsilon_0}\frac{1}{|\vec{x}|}
\end{eqnarray*}
tatsächlich ein hermitescher Operator ist. Zunächst ist der Potential-Anteil $\frac{e^2}{4\pi\epsilon_0}\frac{1}{|\vec{x}|}$ einfach ein reelle Funktion $V(\vec{x})$. Solche können immer als hermitesche Operatoren aufgefasst werden:
\begin{eqnarray*}
\braket{\phi(\vec{x})}{V(\vec{x})\chi(\vec{x})} &=& \int dx^3\ \phi^*(\vec{x})\left(V(\vec{x})\chi(\vec{x})\right)\\
&=& \int dx^3\ V(\vec{x})\left(\phi^*(\vec{x})\chi(\vec{x})\right)\\
&=& \int dx^3\ \left(V(\vec{x})\phi^*(\vec{x})\right)\chi(\vec{x})\\
&=& \braket{V(\vec{x})\phi(\vec{x})}{\chi(\vec{x})}
\end{eqnarray*}
Für den kinetischen Teil bleibt zu zeigen, dass $\Delta$ ein hermitescher Operator ist:
\begin{eqnarray*}
\braket{\phi(\vec{x})}{\Delta\chi(\vec{x})} &=& \int dx^3\ \phi^*(\vec{x})\left(\Delta\chi(\vec{x})\right)\\
&=& \int dx^3\ \phi^*(\vec{x})\left(\nabla\cdot\nabla\chi(\vec{x})\right)\\
&=& \int dx^3\ \nabla\cdot\left(\phi^*(\vec{x})\nabla\chi(\vec{x})\right)-\left(\nabla\phi^*(\vec{x})\right)\cdot\left(\nabla\chi(\vec{x})\right)
\end{eqnarray*}
Das Integral über die Divergenz verschwindet, da normierbare Funktionen im Unendlichen verschwinden.
\begin{eqnarray*}
\braket{\phi(\vec{x})}{\Delta\chi(\vec{x})} &=& -\int dx^3\ \left(\nabla\phi^*(\vec{x})\right)\cdot\left(\nabla\chi(\vec{x})\right)\\
 &=& -\int dx^3\ \nabla\cdot\left(\nabla\phi^*(\vec{x})\right)\left(\chi(\vec{x})\right)-\left(\nabla\cdot\nabla\phi^*(\vec{x})\right)\left(\chi(\vec{x})\right)
\end{eqnarray*}
Erneut verschwindet die Divergenz und wir erhalten schließlich
\begin{eqnarray*}
\braket{\phi(\vec{x})}{\Delta\chi(\vec{x})} &=& \int dx^3\ \left(\nabla\cdot\nabla\phi^*(\vec{x})\right)\left(\chi(\vec{x})\right)\\
 &=& \int dx^3\ \left(\Delta\phi^*(\vec{x})\right)\left(\chi(\vec{x})\right)\\
&=&\braket{\Delta\phi(\vec{x})}{\chi(\vec{x})}
\end{eqnarray*}
\end{subsubsection}

\begin{subsubsection}{Eigenfunktionen zu verschiedenen Eigenwerten eines hermiteschen Operators sind orthogonal}
\label{eigen_normality}
Es gelte für einen hermiteschen Operator $\hat{A}$ und zei Funktionen $\phi$ und $\chi$:
\begin{eqnarray*}
\hat{A}\phi(\vec{x}) &=& a_\phi\phi(\vec{x})\\
\hat{A}\chi(\vec{x}) &=& a_\chi\chi(\vec{x})
\end{eqnarray*}
wobei die beiden reellen $a_\phi$ $a_\chi$ Eigenwerte verschieden seien.
\begin{eqnarray*}
\braket{\phi}{\hat{A}\chi}&=&\braket{\hat{A}\phi}{\chi}\\
\braket{\phi}{a_\phi\chi}&=&\braket{a_\chi\phi}{\chi}\\
a_\phi\braket{\phi}{\chi}&=&a_\chi\braket{\phi}{\chi}\\
\left(a_\phi-a_\chi\right)\braket{\phi}{\chi}&=& 0\\
\braket{\phi}{\chi} &=& 0
\end{eqnarray*}
\end{subsubsection}

\begin{subsubsection}{Aufenthalts-Wahrscheinlichkeit und stationäre Ladungsdichte}
Für ein bestimmtes Energie-Niveau $E_n$ gilt
\begin{eqnarray*}
\phi_n(t,\vec{x}) &=& e^{-i\omega_n t}\sum_{lm} a_{nlm}\phi_{nlm}(\vec{x})
\end{eqnarray*}
Multipliziert man dies mit der komplex konjugierten Welle $\phi^*_n(t,\vec{x})$ erhalten wir eine zeitunabhängige, reelle und positive Funktion $\rho(\vec{x})$, eben die Amplitude einer stehenden Welle:
\begin{eqnarray*}
\rho(\vec{x})&=&\phi^*_n(t,\vec{x})\phi_n(t,\vec{x})\phi_n(t,\vec{x})\phi_n(t,\vec{x})\\
&=& \left(e^{i\omega_n t}\sum_{\lambda\mu} a^*_{n\lambda\mu}\phi^*_{n\lambda\mu}(\vec{x})\right)\left(e^{-i\omega_n t}\sum_{lm} a_{nlm}\phi_{nlm}(\vec{x})\right)\\
&=& \left(\sum_{\lambda\mu} a^*_{n\lambda\mu}\phi^*_{n\lambda\mu}(\vec{x})\right)\left(\sum_{lm} a_{nlm}\phi_{nlm}(\vec{x})\right)
\end{eqnarray*}
Normiert man diese über den Raum, sodass
\begin{eqnarray*}
\int dx^3 \rho(\vec{x}) &=& 1 = \sum_{lm} a^*_{nlm}a_{nlm}=\sum_{lm} |a_{nlm}|^2
\end{eqnarray*}
wird, kann $\rho(\vec{x})$ als räumlich Wahrscheinlichkeits-Verteilung oder Aufenthalts-Wahrscheinlichkeit interpretiert werden.
Multipliziert man diese mit der Elektronen-Ladung $-e$, kann dies als stationäre Ladungsdichte verstanden werden. Deshalb findet keine Abstrahlung von elektromagnetischer Energie statt.
\end{subsubsection}

\begin{subsubsection}{Erwartungswerte und deren Zeitableitung}
Zu einem beliebigen hermiteschen Operator nennen wir
\begin{eqnarray*}
\bra{\phi}\hat{A}\ket{\phi} &=& \int dx^3 \phi^*(\vec{x})\hat{A}\phi(\vec{x})
\end{eqnarray*}
den Erwartungswert von $\hat{A}$. Aus dr zeitabhängigen Schrödinger-Gleichung:
\begin{eqnarray*}
i\hbar \dot{\phi}(t,\vec{x}) &=& \hat{H}\phi(t,\vec{x})\\
\dot{\phi}(t,\vec{x}) &=& \frac{-i}{\hbar}\hat{H}\phi(t,\vec{x})
\end{eqnarray*}
erhalten wir die Zeitableitung der Wellenfunktion, und somit des Erwartungswertes:
\begin{eqnarray*}
\frac{\partial}{\partial t}\bra{\phi}\hat{A}\ket{\phi} &=& \int dx^3 \dot{\phi}^*(\vec{x})\hat{A}\phi(\vec{x})+\phi^*(\vec{x})\hat{A}\dot{\phi}(\vec{x})\\
 &=& \int dx^3 \left(\frac{-i}{\hbar}\hat{H}\phi(t,\vec{x})\right)^*\hat{A}\phi(\vec{x})
+\phi^*(\vec{x})\hat{A}\left(\frac{-i}{\hbar}\hat{H}\phi(t,\vec{x})\right)\\
 &=& \frac{i}{\hbar}\int dx^3 \left(\hat{H}\phi(t,\vec{x})\right)^*\hat{A}\phi(\vec{x})
-\phi^*(\vec{x})\hat{A}\left(\hat{H}\phi(t,\vec{x})\right)\\
&=&\frac{i}{\hbar}\bra{\hat{H}\phi(t,\vec{x})}\hat{A}\ket{\phi(t,\vec{x})}
-\frac{i}{\hbar}\bra{\phi(t,\vec{x})}\hat{A}\ket{\hat{H}\phi(t,\vec{x})}\\
&=&\frac{i}{\hbar}\bra{\phi(t,\vec{x})}\hat{H}\hat{A}\ket{\phi(t,\vec{x})}
-\frac{i}{\hbar}\bra{\phi(t,\vec{x})}\hat{A}\hat{H}\ket{\phi(t,\vec{x})}\\
&=&\frac{i}{\hbar}\bra{\phi(t,\vec{x})}[\hat{H},\hat{A}]\ket{\phi(t,\vec{x})}
\end{eqnarray*}
Formal können wir schreiben:
\begin{eqnarray*}
\frac{\partial}{\partial t}\hat{A} &=&\frac{i}{\hbar}[\hat{H},\hat{A}]
\end{eqnarray*}
\end{subsubsection}

\begin{subsubsection}{Korrespondenz-Prinzip}
Die Zeitableitung lässt sich auch auf so einfache Operatoren wie $\hat{\vec{x}}$
\begin{eqnarray*}
\bra{\phi}\hat{\vec{x}}\ket{\phi} &=& \int dx^3 \phi^*(\vec{x})\vec{x}\phi(\vec{x})
\end{eqnarray*}
anwenden:
\begin{eqnarray*}
\frac{\partial}{\partial t}\hat{\vec{x}} &=&\frac{i}{\hbar}[\hat{H},\hat{\vec{x}}]\cr
[\hat{H},\hat{\vec{x}}]\phi(\vec{x}) &=& -\frac{\hbar^2\Delta}{2m}\vec{x}\phi(\vec{x})
+\vec{x}\frac{\hbar^2\Delta}{2m}\phi(\vec{x})\\
&=& \frac{\hbar^2}{2m}\left(-\Delta\vec{x}\phi(\vec{x})+\vec{x}\Delta\phi(\vec{x})\right)\\
\bra{\phi}\Delta\vec{x}-\vec{x}\Delta\ket{\phi}&=&\int dx^3\ \phi^*(\vec{x})\left(\Delta\vec{x}-\vec{x}\Delta\right)\phi(\vec{x})\\
\left(\Delta\vec{x}\phi(\vec{x})\right)_k &=& \sum_l\partial_l\partial_lx_k\phi(\vec{x})\\
&=& \sum_l\partial_l\delta_{lk}\phi(\vec{x})+\sum_l\partial_lx_k\partial_l\phi(\vec{x})\\
&=& \partial_k\phi(\vec{x})+\sum_l\delta_{lk}\partial_l\phi(\vec{x})+\sum_lx_k\partial_l\partial_l\phi(\vec{x})\\
&=& 2\partial_k\phi(\vec{x})+\sum_lx_k\partial_l\partial_l\phi(\vec{x})\\
\Delta\vec{x}\phi(\vec{x}) &=& 2\vec{\nabla}\phi(\vec{x})+\vec{x}\Delta\phi(\vec{x})\\
\left(\Delta\vec{x}-\vec{x}\Delta\right)\phi(\vec{x}) &=& 2\vec{\nabla}\phi(\vec{x})\\
\bra{\phi}\Delta\vec{x}-\vec{x}\Delta\ket{\phi}&=&2\bra{\phi}\vec{\nabla}\ket{\phi}\cr
[\hat{H},\hat{\vec{x}}] &=& -\frac{\hbar^2}{m}\vec{\nabla}\\
\frac{\partial}{\partial t}\bra{\phi}\hat{\vec{x}}\ket{\phi} &=& \frac{-i}{\hbar}\frac{\hbar^2}{m}\bra{\phi}\vec{\nabla}\ket{\phi}\\
\dot{\hat{\vec{x}}} &=& \hat{\vec{v}}  = \frac{-i\hbar}{m}\vec{\nabla}\\
\hat{\vec{p}} &=& -i\hbar\vec{\nabla}
\end{eqnarray*}
$\hat{\vec{v}}$ ist der Geschwindigkeits-Operator, $\hat{\vec{p}}$ der Impulsoperator. So wird aus der klassischen Hamilton-Funktion
\begin{eqnarray*}
H(p,x) &=& \frac{\vec{p}^2}{2m}+V(\vec{x})
\end{eqnarray*}
Der Hamilton-Operator
\begin{eqnarray*}
\hat{H} &=& \frac{\hat{\vec{p}}^2}{2m}+V(\vec{x})=-\frac{\hbar^2\Delta}{2m}+V(\vec{x})
\end{eqnarray*}
Nun wenden wir dasselbe Prozedere auf $\hat{\vec{p}}$ an:
\begin{eqnarray*}
\frac{\partial}{\partial t}\hat{\vec{p}} &=&\frac{i}{\hbar}[\hat{H},\hat{\vec{p}}]\\
&=& \frac{i}{\hbar}\left[\frac{\hat{\vec{p}}^2}{2m}+V(\vec{x}),\hat{\vec{p}}\right]\\
&=& \frac{i}{\hbar}\left[V(\vec{x}),\hat{\vec{p}}\right]=\left[V(\vec{x}),\hat{\vec{\nabla}}\right]\\
&=&V(\vec{x})\hat{\vec{\nabla}}-\hat{\vec{\nabla}}V(\vec{x})=-\vec{\nabla}V\\
\frac{\partial}{\partial t}\hat{\vec{p}} &=&-\vec{\nabla}V
\end{eqnarray*}
Wir fassen zusammen:
\begin{eqnarray*}
\dot{\hat{\vec{x}}} &=& \frac{\hat{\vec{p}}}{m}\\
\dot{\hat{\vec{p}}} &=&-\vec{\nabla}V
\end{eqnarray*}
und vergleichen dies mit den Hamilton-Gleichungen
\begin{eqnarray*}
\dot{\vec{x}} &=& \frac{\partial H(\vec{p},\vec{x})}{\partial \vec{p}}\\
\dot{\vec{p}} &=&-\frac{\partial H(\vec{p},\vec{x})}{\partial \vec{x}}
\end{eqnarray*}
Wir erhalten also folgende Korrespondenz:
\begin{eqnarray*}
\dot{\hat{\vec{x}}}=\frac{i}{\hbar}[\hat{H},\hat{\vec{x}}] &\sim&\dot{\vec{x}} =\frac{\partial H(\vec{p},\vec{x})}{\partial \vec{p}}\\
\dot{\hat{\vec{p}}}=-\frac{i}{\hbar}[\hat{H},\hat{\vec{p}}] &\sim&\dot{\vec{p}} =-\frac{\partial H(\vec{p},\vec{x})}{\partial \vec{x}}\\
p = \frac{\partial L}{\partial \dot{x}}(\dot{x},x)&\sim& [p,x] =-i\hbar
\end{eqnarray*}
\end{subsubsection}

\begin{subsubsection}{Heisenbergsche Unschärfe-Relation}
\end{subsubsection}
\end{subsection}


\end{section}

\begin{section}{Emission}
\begin{subsection}{Schwebungen und Strom}
Zurück zur Wellenfunktion \ref{double_wave_function}
\begin{eqnarray*}
\phi(t,\vec{x}) &=& a_1\phi_1(t,\vec{x})+a_2\phi_2(t,\vec{x})
\end{eqnarray*}
mit
\begin{eqnarray*}
\hat{H}\Phi_1(\vec{x}) &=& E_1\Phi_1(\vec{x})\\
\braket{\Phi_1(\vec{x})}{\Phi_1(\vec{x})} &=& 1\\
\hat{H}\Phi_2(\vec{x}) &=& E_2\Phi_2(\vec{x})\\
\braket{\Phi_2(\vec{x})}{\Phi_2(\vec{x})} &=& 1\\
E_1 &\ne& E_2\\
\phi_1(t,\vec{x}) &=& e^{-i\omega_1 t}\Phi_1(\vec{x})\\
\phi_2(t,\vec{x}) &=& e^{-i\omega_2 t}\Phi_2(\vec{x})\\
\omega_1 &=& \frac{E_1}{\hbar}\\
\omega_2 &=& \frac{E_2}{\hbar}
\end{eqnarray*}
Die Funktion
\begin{eqnarray*}
\rho(t,\vec{x}) &=& \phi^*(t,\vec{x})\phi(t,\vec{x})
\end{eqnarray*}
ist reell, positiv und normiert, aber nicht konstant. Multipliziert mit der elektrischen Ladung $-e$ des Elektrons
erhalten wir eine zeitabhängige Ladungsverteilung der Gesamtladung $-e$:
\begin{eqnarray*}
\rho_{EM}(t,\vec{x}) &=& -e\phi^*(t,\vec{x})\phi(t,\vec{x})
\end{eqnarray*}
Um die elektrische Stromdichte zu erhalten, leiten wir nach der Zeit ab (wir lassen vorübergehend den Vorfaktor $-e$ weg):
\begin{eqnarray*}
\dot{\rho}(t,\vec{x}) &=& \dot{\phi}^*(t,\vec{x})\phi(t,\vec{x})+\phi^*(t,\vec{x})\dot{\phi}(t,\vec{x})\\
&=& \left(-\frac{i}{\hbar}\hat{H}\phi(t,\vec{x})\right)^*\phi(t,\vec{x})+\phi^*(t,\vec{x})\left(-\frac{i}{\hbar}\hat{H}\phi(t,\vec{x})\right)
\end{eqnarray*}
Dabei haben wir die Schrödinger-Gleichung
\begin{eqnarray*}
i\hbar\dot{\phi}(t,\vec{x}) &=& \hat{H}\phi(t,\vec{x})\\
\dot{\phi}(t,\vec{x}) &=& -\frac{i}{\hbar}\hat{H}\phi(t,\vec{x})
\end{eqnarray*}
benutzt.
\begin{eqnarray*}
\dot{\rho}(t,\vec{x}) &=& \left(-\frac{i}{\hbar}\hat{H}\phi(t,\vec{x})\right)^*\phi(t,\vec{x})+\phi^*(t,\vec{x})\left(-\frac{i}{\hbar}\hat{H}\phi(t,\vec{x})\right)\\
&=& \frac{i}{\hbar}\left(\left(\hat{H}\phi(t,\vec{x})\right)^*\phi(t,\vec{x})-\phi^*(t,\vec{x})\left(\hat{H}\phi(t,\vec{x})\right)\right)\\
-i\hbar\dot{\rho}(t,\vec{x})&=& \left(\hat{H}\phi(t,\vec{x})\right)^*\phi(t,\vec{x})-\phi^*(t,\vec{x})\left(\hat{H}\phi(t,\vec{x})\right)\\
&=& \left(\hat{H}\phi^*(t,\vec{x})\right)\phi(t,\vec{x})-\phi^*(t,\vec{x})\left(\hat{H}\phi(t,\vec{x})\right)
\end{eqnarray*}
Ist $\phi$ eine Eigen-Funktion mit dem Eigenwert $E$, dann verschwindet der obige Ausdruck:
\begin{eqnarray*}
&&\left(\hat{H}\phi^*(t,\vec{x})\right)\phi(t,\vec{x})-\phi^*(t,\vec{x})\left(\hat{H}\phi(t,\vec{x})\right)\\
&=&\left(E\phi^*(t,\vec{x})\right)\phi(t,\vec{x})-\phi^*(t,\vec{x})\left(E\phi(t,\vec{x})\right)\\
&=&E\left(\left(\phi^*(t,\vec{x})\right)\phi(t,\vec{x})-\phi^*(t,\vec{x})\left(\phi(t,\vec{x})\right)\right)=0
\end{eqnarray*}
Betrachten wir nur den Potential-Anteil $V(\vec{x})=\frac{e^2}{4\pi\epsilon_0}\frac{1}{|\vec{x}|}$, dann verschwindet der Ausdruck ebenso:
\begin{eqnarray*}
&&\left(V(\vec{x})\phi^*(t,\vec{x})\right)\phi(t,\vec{x})-\phi^*(t,\vec{x})\left(V(\vec{x})\phi(t,\vec{x})\right)\\
&=&V(\vec{x})\left(\phi^*(t,\vec{x})\phi(t,\vec{x})-\phi^*(t,\vec{x})\phi(t,\vec{x})\right)=0
\end{eqnarray*}
Wir können also $\hat{H}$ durch $-\frac{\hbar^2}{2m}\Delta$ ersetzen:
\begin{eqnarray*}
-i\hbar\dot{\rho}(t,\vec{x}) &=&-\frac{\hbar^2}{2m}\left(\left(\Delta\phi^*(t,\vec{x})\right)\phi(t,\vec{x})-\phi^*(t,\vec{x})\left(\Delta\phi(t,\vec{x})\right)\right)\\
i\frac{2m}{\hbar}\dot{\rho}(t,\vec{x}) &=&\left(\Delta\phi^*(t,\vec{x})\right)\phi(t,\vec{x})-\phi^*(t,\vec{x})\left(\Delta\phi(t,\vec{x})\right)
\end{eqnarray*}
Wir behandeln die rechte Seite weiter:
\begin{eqnarray*}
&&\left(\Delta\phi^*(t,\vec{x})\right)\phi(t,\vec{x})-\phi^*(t,\vec{x})\left(\Delta\phi(t,\vec{x})\right)\\
&=&\nabla\cdot\left(\left(\nabla\phi^*(t,\vec{x})\right)\phi(t,\vec{x})\right)-\left(\nabla\phi^*(t,\vec{x})\right)\cdot\left(\nabla\phi(t,\vec{x})\right)\\
&&-\nabla\cdot\left(\left(\phi^*(t,\vec{x})\right)\nabla\phi(t,\vec{x})\right)+\left(\nabla\phi^*(t,\vec{x})\right)\cdot\left(\nabla\phi(t,\vec{x})\right)\\
&=&\nabla\cdot\left(\left(\nabla\phi^*(t,\vec{x})\right)\phi(t,\vec{x})-\left(\phi^*(t,\vec{x})\right)\nabla\phi(t,\vec{x})\right)
\end{eqnarray*}
und erhalten schließlich
\begin{eqnarray*}
i\frac{2m}{\hbar}\dot{\rho}(t,\vec{x})&=&\nabla\cdot\left(\left(\nabla\phi^*(t,\vec{x})\right)\phi(t,\vec{x})-\left(\phi^*(t,\vec{x})\right)\nabla\phi(t,\vec{x})\right)\\
\dot{\rho}(t,\vec{x})&=&-i\frac{\hbar}{2m}\nabla\cdot\left(\left(\nabla\phi^*(t,\vec{x})\right)\phi(t,\vec{x})-\left(\phi^*(t,\vec{x})\right)\nabla\phi(t,\vec{x})\right)
\end{eqnarray*}
Definieren wir den Wahrscheinlichkeits-Strom
\begin{eqnarray*}
\vec{j} &=& i\frac{\hbar}{2m}\left(\left(\nabla\phi^*(t,\vec{x})\right)\phi(t,\vec{x})-\left(\phi^*(t,\vec{x})\right)\nabla\phi(t,\vec{x})\right)
\end{eqnarray*}
, dann erhalten wir eine Kontinuitäts-Gleichung für die Wahrscheinlichkeit:
\begin{eqnarray*}
\vec{j} &=& -\nabla\cdot\vec{j}\\
\vec{j}+\nabla\cdot\vec{j} &=& 0
\end{eqnarray*}
Die elektrische Stromdichte ist also
\begin{eqnarray*}
\vec{j}_{EM} &=& -e\vec{j}
\end{eqnarray*}
Weder $\rho_{EM}$ noch $\vec{j}_{EM}$ sind messbare Größen, aber sie können dennoch als Quellen eines elektromagnetischen Feldes betrachtet werden.

Nun zur Wellenfunktion \ref{double_wave_function}
\begin{eqnarray*}
\phi(t,\vec{x}) &=& a_1\phi_1(t,\vec{x})+a_2\phi_2(t,\vec{x})\\
&=& a_1e {-i\omega_1}\Phi_1(\vec{x})+a_2e^{-i\omega_2}\Phi_2(\vec{x})
\end{eqnarray*}
Die Aufenthaltswahrscheinlichkeit ist
\begin{eqnarray*}
\rho(t,\vec{x}) &=& \left(a_1\phi_1(t,\vec{x})+a_2\phi_2(t,\vec{x})\right)^*\left(a_1\phi_1(t,\vec{x})+a_2\phi_2(t,\vec{x})\right)\\
&=&a_1^*a_1\Phi^*_1(\vec{x})\Phi_1(\vec{x})+a_2^*a_2\Phi^*_2(\vec{x})\Phi_2(\vec{x})\\
&&+a_1^*a_2e^{i(\omega_1-\omega_2)t}\Phi^*_1(\vec{x})\Phi_2(\vec{x})+a_2^*a_1e^{-i(\omega_2-\omega_1)t}\Phi^*_2(\vec{x})\Phi_1(\vec{x})
\end{eqnarray*}
Der erste Teil beschreibt eine stationäre Aufenthaltswahrscheinlichkeit, der zweite den oszillierenden.
\begin{eqnarray*}
\rho_{osc}(t,\vec{x}) &=& a_1^*a_2e^{i(\omega_1-\omega_2)t}\Phi^*_1(\vec{x})\Phi_2(\vec{x})+a_2^*a_1e^{-i(\omega_2-\omega_1)t}\Phi^*_2(\vec{x})\Phi_1(\vec{x})\\
&=& 2Re\left(a_1^*a_2e^{i(\omega_1-\omega_2)t}\Phi^*_1(\vec{x})\Phi_2(\vec{x})\right)
\end{eqnarray*}
und analog für den Wahrscheinlichkeits-Strom:
\begin{eqnarray*}
\vec{j}(t,\vec{x}) &=& i\frac{\hbar}{2m}\left(\left(\nabla\phi^*(t,\vec{x})\right)\phi(t,\vec{x})-\left(\phi^*(t,\vec{x})\right)\nabla\phi(t,\vec{x})\right)\\
&=& \frac{\hbar}{m}Re\left(\left(\nabla\phi^*(t,\vec{x})\right)\phi(t,\vec{x})\right)\\
\left(\nabla\phi^*(t,\vec{x})\right)\phi(t,\vec{x}) &=& e^{i(\omega_1-\omega_2)t}\left(\nabla\Phi_1^*(\vec{x})\right)\Phi_2(\vec{x})\\
&&+e^{-i(\omega_1-\omega_2)t}\left(\nabla\Phi_2^*(\vec{x})\right)\Phi_1(\vec{x})\\
\vec{j}_{osc} &=& \frac{\hbar}{m}Re\left(e^{i(\omega_1-\omega_2)t}\left(\nabla\Phi_1^*(\vec{x})\right)\Phi_2(\vec{x})\right)\\
&&+\frac{\hbar}{m}Re\left(e^{-i(\omega_1-\omega_2)t}\left(\nabla\Phi_2^*(\vec{x})\right)\Phi_1(\vec{x})\right)
\end{eqnarray*}
\end{subsection}

\begin{subsection}{Dipolstrahlung}
bjkuigg
hhhh

\end{subsection}
\end{section}

\ifx\makebook\undefined
	\end{document}
\else
	\end{chapter}
\fi
